\documentclass[11pt,a4paper]{article}

\usepackage[utf8]{inputenc} % allow utf-8 input
\usepackage[T1]{fontenc}    % use 8-bit T1 fonts
\usepackage{hyperref}       % hyperlinks
\usepackage{url}            % simple URL typesetting
\usepackage{booktabs}       % professional-quality tables
\usepackage{amsfonts}       % blackboard math symbols
\usepackage{nicefrac}       % compact symbols for 1/2, etc.
\usepackage{microtype}      % microtypography
\usepackage{xcolor}         % colors
\usepackage{graphicx}       % graphics
\usepackage[normalem]{ulem} % strikethrough

\begin{document}

\begin{table}[h]
\begin{tabular}{|l|}
\hline
SATURATED HYDRAULIC CONDUCTIVITY: INVERSE AUGERHOLE METHOD \\ \hline
\end{tabular}
\end{table}

\begin{table}[h]
\begin{tabular}{|l|l|l|l|l|l|l|l|l|}
\hline
NO: & 4003 &   & TEXT: & sal & SOIL: & c &  &  \\ \hline
DEPTH(D): & 30 & cm & Radius(r): & 3 & cm & L(s+f) & 46.9 & cm \\ \hline
\end{tabular}
\end{table}

\begin{table}[h]
\begin{tabular}{|l|l|l|l|l|}
\hline
i & t(s) & h'(cm) & h(cm) & log(h+r/2) \\ \hline
1 & 0 &  & -16.9 & ERR \\ \hline
2 & 30 &  & -16.9 & ERR \\ \hline
3 & 60 &  & -16.9 & ERR \\ \hline
4 & 90 &  & -16.9 & ERR \\ \hline
5 & 120 &  & -16.9 & ERR \\ \hline
6 & 150 &  & -16.9 & ERR \\ \hline
7 & 180 &  & -16.9 & ERR \\ \hline
8 & 210 &  & -16.9 & ERR \\ \hline
9 & 240 &  & -16.9 & ERR \\ \hline
10 & 270 &  & -16.9 & ERR \\ \hline
11 & 300 &  & -16.9 & ERR \\ \hline
12 & 330 &  & -16.9 & ERR \\ \hline
13 & 360 &  & -16.9 & ERR \\ \hline
14 & 390 &  & -16.9 & ERR \\ \hline
15 & 420 &  & -16.9 & ERR \\ \hline
16 & 450 &  & -16.9 & ERR \\ \hline
17 & 480 &  & -16.9 & ERR \\ \hline
18 & 510 &  & -16.9 & ERR \\ \hline
19 & 540 &  & -16.9 & ERR \\ \hline
20 & 570 &  & -16.9 & ERR \\ \hline
21 & 600 &  & -16.9 & ERR \\ \hline
22 & 630 &  & -16.9 & ERR \\ \hline
23 & 660 &  & -16.9 & ERR \\ \hline
24 & 690 &  & -16.9 & ERR \\ \hline
25 & 720 &  & -16.9 & ERR \\ \hline
26 & 750 &  & -16.9 & ERR \\ \hline
27 & 780 &  & -16.9 & ERR \\ \hline
28 & 810 &  & -16.9 & ERR \\ \hline
29 & 840 &  & -16.9 & ERR \\ \hline
30 & 870 &  & -16.9 & ERR \\ \hline
31 & 900 &  & -16.9 & ERR \\ \hline
32 & 930 &  & -16.9 & ERR \\ \hline
\end{tabular}
\end{table}

\begin{table}[h]
\begin{tabular}{|l|}
\hline
Regression Output: \\ \hline
\end{tabular}
\end{table}

\begin{table}[h]
\begin{tabular}{|l|}
\hline
Constant \\ \hline
Std Err of Y Est \\ \hline
R Squared \\ \hline
No. of Observations \\ \hline
Degrees of Freedom \\ \hline
\end{tabular}
\end{table}

\begin{table}[h]
\begin{tabular}{|l|}
\hline
3.999287172371125 \\ \hline
0.000000000000001012679207245471 \\ \hline
0.9999999999996305 \\ \hline
32 \\ \hline
30 \\ \hline
\end{tabular}
\end{table}

\begin{table}[h]
\begin{tabular}{|l|}
\hline
X Coefficient(s) \\ \hline
Std Err of Coef. \\ \hline
\end{tabular}
\end{table}

\begin{table}[h]
\begin{tabular}{|l|}
\hline
-0.000000000000000003398892096748884 \\ \hline
0.0000000000000000006462915165271856 \\ \hline
\end{tabular}
\end{table}

\begin{table}[h]
\begin{tabular}{|l|l|l|}
\hline
Ks: & 0.000000000000010131417561989074 & (m/day) \\ \hline
Std.Err: & 0.000000000000001926465752464235 & (m/day) \\ \hline
\end{tabular}
\end{table}

\end{document}